\documentclass[11pt]{amsart}

\usepackage[T1]{fontenc}
\usepackage[utf8]{inputenc}
\IfFileExists{lmodern.sty}{\usepackage{lmodern}}{\PassOptionsToPackage{expansion=false}{microtype}}
\usepackage{amsmath,amssymb,amsthm,mathtools}
\usepackage{enumitem}
\usepackage[nopatch=footnote]{microtype}
\usepackage{xcolor}
\usepackage[colorlinks=true,linkcolor=blue,citecolor=blue,urlcolor=blue]{hyperref}
\hypersetup{
  pdftitle={Coordinate-Block Factorization under Single-Sandwich Division and Applications to Property (TT)/T},
  pdfauthor={GPT-5.6 Sol},
  pdfsubject={quasi-cocycles, elementary groups, Leavitt algebras},
  pdfkeywords={property (TT)/T, quasi-cocycles, elementary groups, Leavitt algebras, bounded generation}
}

\newif\ifanonymous\anonymousfalse
\ifanonymous
  \newcommand{\leanrepo}{[repository link removed for anonymous review]}
\else
  \newcommand{\leanrepo}{https://github.com/SauersML/group-approximation}
\fi
% Machine-readable theorem correspondence, intentionally invisible in print.
\newcommand{\leanverified}[2]{}

\setlist[enumerate]{leftmargin=2.35em}
\setlist[itemize]{leftmargin=2.1em}
\allowdisplaybreaks
\raggedbottom
\emergencystretch=1.5em
\extrafloats{100}

\newtheorem{theorem}{Theorem}[section]
\newtheorem{proposition}[theorem]{Proposition}
\newtheorem{lemma}[theorem]{Lemma}
\newtheorem{corollary}[theorem]{Corollary}
\theoremstyle{definition}
\newtheorem{definition}[theorem]{Definition}
\newtheorem{remark}[theorem]{Remark}

\theoremstyle{plain}
\newtheorem{mainthm}{Theorem}
\renewcommand{\themainthm}{\Alph{mainthm}}

\newcommand{\F}{\mathbb F}
\newcommand{\N}{\mathbb N}
\newcommand{\C}{\mathbb C}
\newcommand{\GL}{\operatorname{GL}}
\newcommand{\EL}{\operatorname{E}}
\newcommand{\diag}{\operatorname{diag}}
\newcommand{\Fix}{\operatorname{Fix}}
\newcommand{\Def}{\operatorname{def}}
\newcommand{\id}{\mathrm{id}}
\newcommand{\norm}[1]{\left\lVert#1\right\rVert}
\newcommand{\abs}[1]{\left\lvert#1\right\rvert}

\title[Coordinate-block factorization and \texorpdfstring{\((TT)/T\)}{(TT)/T}]
{Coordinate-Block Factorization under Single-Sandwich Division
and Applications to Property \texorpdfstring{\((TT)/T\)}{(TT)/T}}
\author{GPT-5.6 Sol}
\date{August 2026}
\subjclass[2020]{20J06, 20F65, 16S50, 16S88, 19B99}
\keywords{property (TT)/T, quasi-cocycles, elementary groups,
Leavitt algebras, purely infinite rings, bounded generation}

\begin{document}

\begin{abstract}
We prove a coordinate-block factorization theorem for general linear
groups over rings satisfying the single-sandwich condition
\[
  a\ne0\quad\Longrightarrow\quad
  xay=1\ \text{for some }x,y.
\]
Every element of \(\GL_n(R)\) is a product of at most \(2n+6\)
factors, each an elementary transvection or an element of one fixed
coordinate copy of \(\GL_{n-1}(R)\).  Thus \(\GL_n(R)\) has width at
most \(2n+6\) relative to the elementary root subgroups and that
fixed coordinate block.

The algebraic factorization is the main input developed here.  We
combine it with Mimura's relative quasi-cocycle rigidity framework
for elementary groups over finite free associative
\(\F_2\)-algebras.  If \(R\) is a nontrivial finite-type
\(\F_2\)-algebra carrying a binary Leavitt family, satisfying the
single-sandwich condition, and satisfying
\(\diag(u,1)\in\EL_2(R)\) for every \(u\in R^\times\), then
\(\EL_n(R)\) has property \((TT)/T\) for every \(n\ge3\).
Consequently,
\[
  \EL_n\bigl(L_{\F_2}(1,2)\bigr)
  \quad\text{has property }(TT)/T
  \qquad(n\ge3).
\]
An independent quantitative reconstruction of the relative root
estimate makes the analytic input explicit.  Together with OpenAI's
previously
established nonsoficity theorem for the binary Leavitt unit
group~\cite[Theorem~1.1, Chapter~3]{OpenAITenProofs}, the result gives explicit
nonsofic groups with property \((TT)/T\).
\end{abstract}

\maketitle

\section{Introduction}\label{sec:introduction}

Property \((TT)/T\), introduced by Mimura, asks that Hilbert-space
quasi-cocycles be bounded after the invariant-vector part of the
representation has been removed.  It is stronger than Kazhdan's
property~\((T)\) and is designed to control bounded cohomology and
homomorphisms into groups admitting unbounded
quasi-cocycles~\cite{MimuraTT}.  Ershov and Jaikin-Zapirain proved property~\((T)\)
for elementary groups over arbitrary finitely generated associative
rings~\cite{ErshovJaikin}.  The corresponding quasi-cocycle problem
is subtler: the bounded error in the cocycle identity accumulates
under a long word, so finite generation alone gives no global bound.

The present paper supplies a globalization mechanism for a class of
purely infinite rings.  Its algebraic input is an explicit
matrix factorization.  The hypothesis is the following strong
two-sided division property:
\begin{equation}\label{eq:sandwich}
  a\ne0\quad\Longrightarrow\quad
  \exists x,y\in R,\qquad xay=1.
\end{equation}
We call~\eqref{eq:sandwich} the single-sandwich property.  Ara, Goodearl, and
Pardo proved that, for a nondivision simple ring, this is equivalent
to pure infiniteness~\cite[Theorem~1.6]{AraGoodearlPardo}.

\begin{mainthm}[Coordinate-block factorization]\label{thm:A}
Let \(R\) be a nontrivial ring satisfying~\eqref{eq:sandwich}, let
\(n\ge2\), and fix a coordinate \(j<n\).  Every \(A\in\GL_n(R)\)
admits a factorization of length at most \(2n+6\) in which every
factor is either an elementary transvection or belongs to the
coordinate copy
\[
  \GL_{n-1}^{(j)}(R)
  =\{B\in\GL_n(R):Be_j=e_j,\ e_j^{\,t}B=e_j^{\,t}\}.
\]
\leanverified{PropertyTT/PaperStatements}{PropertyTTPaper.coordinateBlock_factorization}
\end{mainthm}

Theorem~\ref{thm:A} is a relative bounded-width statement rather
than bounded elementary generation: the bounded factors come from
the elementary root subgroups together with one fixed coordinate
block.  This is exactly the form required by the quasi-cocycle
argument.  By the characterization just cited, the conclusion applies
in particular to every unital purely infinite simple ring.

For the synthesis theorem, a \emph{binary Leavitt family} in \(R\)
means elements \(x_0,x_1,y_0,y_1\) satisfying
\[
  y_i x_j=\delta_{ij},\qquad x_0y_0+x_1y_1=1.
\]
It supplies the matrix self-similarity used to move between ranks.
We say that an \(\F_2\)-algebra is \emph{finite type} when it is
generated by a finite set as an \(\F_2\)-algebra, equivalently when
it is a quotient of \(\F_2\langle X\rangle\) for some finite set
\(X\).

\begin{mainthm}[Rigidity synthesis]\label{thm:B}
Let \(R\) be a nontrivial finite-type \(\F_2\)-algebra.  Assume that
\begin{enumerate}
\item \(R\) contains a binary Leavitt family;
\item \(R\) satisfies the single-sandwich
  property~\eqref{eq:sandwich};
\item \(\diag(u,1)\in\EL_2(R)\) for every \(u\in R^\times\).
\end{enumerate}
Then \(\EL_n(R)\) has property \((TT)/T\) for every \(n\ge3\).
\leanverified{PropertyTT/PaperStatements}{PropertyTTPaper.finiteType_elementaryGroup_hasTTmodT}
\end{mainthm}

The assumptions play separate roles: finite type presents the
ring as a quotient of a finite free algebra, the Leavitt family moves
between matrix ranks, the single-sandwich property supplies bounded
matrix reduction, and the diagonal condition identifies elementary
and general linear groups.

\begin{mainthm}[Binary Leavitt algebra]\label{thm:C}
For every \(n\ge3\),
\[
  \EL_n\bigl(L_{\F_2}(1,2)\bigr)
  \quad\text{has property }(TT)/T.
\]
\leanverified{PropertyTT/PaperStatements}{PropertyTTPaper.binaryLeavitt_elementaryGroup_hasTTmodT}
\end{mainthm}

Combining Theorem~\ref{thm:C} with OpenAI's theorem that
\(L_{\F_2}(1,2)^\times\) is nonsofic~\cite[Theorem~1.1 of
Chapter~3]{OpenAITenProofs} gives the following consequence.

\begin{corollary}\label{cor:nonsofic}
For every \(n\ge3\), the group
\(\EL_n(L_{\F_2}(1,2))\) is nonsofic and has property \((TT)/T\).
\leanverified{PropertyTT/NonsoficCorollary}{PropertyTTPaper.binaryLeavitt_elementaryGroup_hasTTmodT_and_not_isSofic}
\end{corollary}

\begin{proof}
Theorem~\ref{thm:C} supplies property \((TT)/T\).  For
nonsoficity, OpenAI's theorem gives that \(L^\times\) is
nonsofic~\cite[Theorem~1.1 of Chapter~3]{OpenAITenProofs}.  The map
\[
  L^\times\longrightarrow\GL_n(L),
  \qquad u\longmapsto\diag(u,1,\ldots,1),
\]
is injective.  Proposition~\ref{prop:diag}, followed by the standard
block inclusion from rank two to rank \(n\), shows that its image
lies in \(\EL_n(L)\).  Since every subgroup of a sofic group is
sofic, \(\EL_n(L)\) cannot be sofic.
\end{proof}

\subsection*{Relation to previous work}

The argument combines several established theories.  Mimura
introduced property \((TT)/T\), developed the relative-to-global
strategy for quasi-cocycles, and proved the relative results for
universal-lattice settings used below~\cite{MimuraTT,MimuraThesis}.
Ershov and Jaikin-Zapirain proved property~\((T)\) for elementary
groups over finitely generated associative rings~\cite{ErshovJaikin}.
On the ring-theoretic side, Ara, Goodearl, and Pardo established the
GE and unit-group \(K_1\) structure of purely infinite simple
rings~\cite{AraGoodearlPardo}, while Ara, Brustenga, and Corti\~nas
computed the stable algebraic \(K\)-theory of Leavitt path
algebras~\cite{AraBrustengaCortinas}.

Within this framework, the algebraic input developed here is the fixed-coordinate
factorization of Theorem~\ref{thm:A}, including its explicit width
\(2n+6\).  The rigidity theorem arises by using that factorization to
globalize relative quasi-cocycle bounds.  Section~\ref{sec:quasicocycles}
gives a quantitative proof of the root estimate, and
Section~\ref{sec:leavitt} verifies the remaining structural
hypotheses for the binary Leavitt algebra.

\subsection*{Organization}

Section~\ref{sec:quasicocycles} develops the relative root estimate.
Section~\ref{sec:factorization} proves the five-move pivot and the
coordinate-block factorization.  Section~\ref{sec:globalization}
combines the analytic and algebraic inputs to prove
Theorem~\ref{thm:B}, and Section~\ref{sec:leavitt} verifies its
hypotheses for the binary Leavitt algebra.  The final section
discusses the natural extensions suggested by the proof.

\section{Quasi-cocycles and relative control}\label{sec:quasicocycles}

Let \(G\) be a group and let
\(\pi:G\to\mathcal U(\mathcal H)\) be a unitary representation on a
complex Hilbert space.  A map \(b:G\to\mathcal H\) is a
\emph{quasi-cocycle of defect at most \(D\)} if
\begin{equation}\label{eq:quasicocycle}
  \norm{b(gh)-b(g)-\pi(g)b(h)}\le D
  \qquad(g,h\in G).
\end{equation}
An exact cocycle has \(D=0\); it is the translation part of an
affine isometric action.  A quasi-cocycle permits uniformly bounded
failure of the affine action law.

\begin{definition}
A group \(G\) has property \((TT)/T\) if, for every unitary
representation \(\pi\) with no nonzero invariant vector, every
Hilbert-space quasi-cocycle into \(\pi\) is bounded on \(G\).
For a subgroup \(H\le G\), the pair \((G,H)\) has relative property
\((TT)/T\) if the same hypotheses imply boundedness on \(H\).  We
write relative \((TT)\) when the conclusion holds without excluding
invariant vectors.
\end{definition}

\paragraph{Mimura's relative input.}
Let \(X\) be finite and \(A=\F_2\langle X\rangle\).  The pair
\(
  (\EL_4(A),X_{03})
\)
has relative property \((TT)/T\).  Indeed,
Mimura proves relative \((TT)/T\) for the noncommutative universal
lattice pair in~\cite[p.~154, Proposition~9.2.8]{MimuraThesis}:
\[
  \EL_m(\mathbb Z\langle X\rangle)
  \ \ge\ \EL_{m-1}(\mathbb Z\langle X\rangle),
  \qquad m\ge3
\].
Relative \((TT)/T\)
descends through a surjective homomorphism: pull a representation
and quasi-cocycle back along the quotient map and then apply the
relative bound.  Reduce coefficients modulo \(2\), take \(m=4\),
and conjugate the standard rank-three block so that it contains
\(X_{03}\).  Boundedness on that block implies boundedness on the
root subgroup.

The local input used below is slightly stronger than a relative
\((TT)/T\) statement because no hypothesis is imposed on invariant
vectors.

We first isolate the quantitative rank-three estimate used in the
direct proof.  For distinct \(i,j,k\), write
\(H_{ijk}=\langle X_{ik},X_{jk}\rangle\le\EL_3(A)\), and identify
the elements of \(X\) with the corresponding free generators of
\(A\).

\begin{lemma}[Localized plane estimate]\label{lem:localized-plane}
Let \(\sigma:\EL_3(A)\to\mathcal U(\mathcal H)\) be a unitary
representation and let \(z\in\mathcal H\).  Suppose \(B\ge0\) and
the displacement of \(z\) is at most \(B\) under
\[
 x_{ik}(1),\ x_{jk}(1),\ x_{ij}(1),\ x_{ji}(1),
 \quad x_{ij}(s),\ x_{ji}(s)\quad(s\in X).
\]
Then, for every \(g\in H_{ijk}\),
\[
  \norm{\sigma(g)z-z}
  \le 2(6\abs X+6)(B+1).
\]
\leanverified{PropertyTT/PaperStatements}{PropertyTTPaper.freeCharacteristicTwo_localizedPlaneEstimate}
\end{lemma}

\begin{proof}
Put \(\delta=B+1\), let \(P\) be the orthogonal projection onto
\(\mathcal H^{H_{ijk}}\), and set \(w=z-Pz\).  The adjacent root
groups \(X_{ij}\) and \(X_{ji}\) normalize \(H_{ijk}\), while
\(X_{ik},X_{jk}\le H_{ijk}\).  Orthogonal projection therefore
passes each of the displayed displacement bounds from \(z\) to
\(w\), with the strict bound \(\delta\).

For completeness, we recall the energy estimate behind this step.
Filter the free algebra by word length and decompose the finite
elementary abelian plane into its characters.  The shears
\(x_{ij}(1),x_{ji}(1),x_{ij}(s),x_{ji}(s)\) move the character mass
between the four leading-letter regions.  Summing the resulting
parallelogram identities and then letting the filtration degree tend
to infinity gives
\begin{align*}
 \norm w^2 \le{}& d_{ik}^2+d_{jk}^2
 +\frac32\sum_{s\in X}
   2\norm w\bigl(d_{ij,s}+d_{ji,s}\bigr)\\
 &+2\norm w\bigl(d_{ij,1}+d_{ji,1}\bigr),
\end{align*}
where each \(d\) is the corresponding displacement of \(w\).
The boundary terms vanish because the word-length filtration is
exhaustive.  Hence
\[
  \norm w^2
  \le 2\delta^2+(6\abs X+4)\norm w\,\delta.
\]
If \(\norm w\le\delta\), the desired bound on \(w\) is immediate.
Otherwise \(\delta\le\norm w\), and the last inequality gives
\(\norm w\le(6\abs X+6)\delta\).  Finally, every \(g\in H_{ijk}\)
fixes \(Pz\), so
\(\norm{\sigma(g)z-z}\le2\norm w\).
\end{proof}

\begin{theorem}[Root quasi-cocycle estimate]\label{thm:root}
Let \(X\) be finite, put \(A=\F_2\langle X\rangle\), and let
\[
  G=\EL_4(A),\qquad
  X_{03}=\{I+aE_{03}:a\in A\}.
\]
For every unitary representation \(\pi\), every \(D\ge0\), and every
quasi-cocycle \(b:G\to\mathcal H\) of defect at most \(D\), the set
\(b(X_{03})\) is bounded.
\leanverified{PropertyTT/PaperStatements}{PropertyTTPaper.freeCharacteristicTwo_root_hasRelativeTT}
\end{theorem}

\begin{proof}
Let \(N\) be the last-column subgroup generated by
\(X_{03},X_{13},X_{23}\).  Decompose
\[
  \mathcal H=\mathcal H^N\oplus(\mathcal H^N)^\perp,
  \qquad b(r)=p+w
\]
for \(r=I+aE_{03}\).  Since the coefficient ring has
characteristic two, \(r^2=1\).
First, \(\norm{b(1)}\le D\), by applying~\eqref{eq:quasicocycle}
to \(1\cdot1\).  Apply the same inequality to \(r^2\) and project
orthogonally to \(\mathcal H^N\).  Since \(r\in N\), it acts as the
identity on this subspace.  If \(q\) denotes the projection of
\(b(1)\), contractivity of the projection gives
\[
  \norm{q-p-p}\le D,
  \qquad \norm q\le D.
\]
Thus \(\norm{2p}\le2D\), and hence
\begin{equation}\label{eq:fixed}
  \norm p\le D.
\end{equation}

For the moving component, let \(Q_X\subseteq\EL_3(A)\) consist of
all elementary roots with coefficient \(1\) or one of the free
generators.  Embed it on the coordinates \(\{1,2,3\}\), and also on
the coordinates \(\{0,1,2\}\) using the transpose--word-reversal
symmetry.  Put
\[
 M=\sum_{s\in Q_X}
 \bigl(\norm{b(\iota_{123}(s))}+\norm{b(\iota_{012}(s))}\bigr),
 \quad B=2M+2D,
\]
\[
 K=2(6\abs X+6)(B+1).
\]
Only the following members of these two finite embedded sets are
used:
\[
 \begin{array}{c}
 x_{13}(1),\ x_{23}(1),\ x_{12}(t),\ x_{21}(t),\\
 x_{01}(1),\ x_{02}(1),\ x_{21}(t),\ x_{12}(t),
 \end{array}
 \qquad t\in\{1\}\cup X.
\]
Here the first row comes from the \(123\)-embedding and the second
from the \(012\)-embedding.  Every displayed element commutes with
\(r=x_{03}(a)\) and normalizes \(N\).  Comparing the
quasi-cocycle equations for \(sr\) and \(rs\), then projecting to
the moving subspace, gives for each displayed control element
\[
  \norm{\pi(s)w-w}
  \le \norm{\pi(s)b(r)-b(r)}
  \le 2\norm{b(s)}+2D\le B.
\]
Lemma~\ref{lem:localized-plane} propagates these finitely many
inequalities to the complete root planes in the two embeddings.
In particular, for every \(c\in A\),
\[
 \norm{\pi(x_{13}(c))w-w},\quad
 \norm{\pi(x_{23}(c))w-w},\quad
 \norm{\pi(x_{01}(c))w-w}
 \le K.
\]
The Steinberg
relation
\[
  [\,I+cE_{01},\,I+E_{13}\,]=I+cE_{03}
\]
and the unitary commutator estimate give displacement at most
\(4K\) under \(x_{03}(c)\).  Every element of \(N\) is a product of
one element from each of \(X_{03},X_{13},X_{23}\), so its
displacement of \(w\) is at most \(6K\).  Since
\(w\perp\mathcal H^N\), this already bounds \(w\).  Indeed, the
closed convex hull of the \(N\)-orbit of \(w\) has a unique vector
of minimal norm; that vector is \(N\)-fixed.  It is therefore zero,
whereas the whole orbit, and hence its closed convex hull, lies in
the closed ball of radius \(6K\) about \(w\).  Consequently
\(\norm w\le6K\).  Together with~\eqref{eq:fixed}, this gives
\[
  \norm{b(I+aE_{03})}
  \le
  D+12(6\abs X+6)(2M+2D+1),
\]
uniformly in \(a\in A\).
\end{proof}

\begin{remark}[A qualitative proof from Mimura's theorem]\label{rem:root-literature}
For qualitative boundedness, Mimura's relative theorem gives a
shorter proof of Theorem~\ref{thm:root}.  Split
\(\mathcal H=\mathcal H^G\oplus(\mathcal H^G)^\perp\).  Mimura's
relative theorem bounds the second projection on \(X_{03}\).  On
the invariant summand the action is trivial, while \(r^2=1\) for
\(r\in X_{03}\); the quasi-cocycle identities for \(1\) and \(r^2\)
then give \(\norm{b_0(r)}\le D\).  The direct proof above complements
this argument by recording an explicit bound and the quantitative
rank-three propagation behind it.
\end{remark}

\section{Coordinate-block factorization}\label{sec:factorization}

We now prove the main algebraic input.  Write
\[
  e_{ij}(a)=I+aE_{ij}\qquad(i\ne j)
\]
for an elementary transvection.

\begin{lemma}[Five-move pivot]\label{lem:pivot}
Let \(R\) satisfy~\eqref{eq:sandwich}, let \(A\in\GL_n(R)\), and
fix distinct coordinates \(i,j\).  At most two elementary row
operations and at most three elementary column operations produce a
matrix \(B\) satisfying
\[
  B_{jj}=1,\qquad B_{ij}=B_{ji}=0.
\]
\leanverified{PropertyTT/PaperStatements}{PropertyTTPaper.fiveMove_pivot}
\end{lemma}

\begin{proof}
We give the calculation, since it is the source of the uniform
constant.  Choose two distinct indices \(i,j\).  First arrange that
the entry \(a=A_{ii}\) is nonzero.  If it is zero, invertibility of
\(A\) gives \(k\) with \(A_{ik}\ne0\); right multiplication by
\(e_{ki}(1)\) replaces \(A_{ii}\) by \(A_{ik}\).  This costs at most
one preliminary column operation.

Write \(b=A_{ij}\), and choose \(p,q\in R\) with \(paq=1\).  Put
\[
  \rho=q(1-pb),\qquad
  F_1=e_{ij}(\rho),\qquad A_1=AF_1,\qquad b_1=(A_1)_{ij}.
\]
Then
\[
  b_1=a\rho+b,\qquad
  pb_1=pa\,q(1-pb)+pb=1.
\]
Set
\[
  d_1=(A_1)_{jj},\qquad
  w=(1-d_1)p,\qquad
  E_1=e_{ji}(w),\qquad A_2=E_1A_1.
\]
The \((j,j)\)-entry is now
\[
  (A_2)_{jj}=d_1+wb_1
  =d_1+(1-d_1)pb_1=1.
\]
Next set
\[
  E_2=e_{ij}(-b_1),\qquad A_3=E_2A_2.
\]
This clears \((i,j)\) without changing the new \(1\) at \((j,j)\).
Finally, with \(c_3=(A_3)_{ji}\), put
\[
  F_2=e_{ji}(-c_3),\qquad A_4=A_3F_2.
\]
This clears \((j,i)\), again preserving \((j,j)=1\) and
\((i,j)=0\).  Thus the four main moves use two row and two column
transvections; including the possible preliminary column addition
gives the asserted bounds \(2\) and \(3\).  No stable-range or
commutativity hypothesis is used.
\end{proof}

\begin{proof}[Proof of Theorem~\ref{thm:A}]
Choose \(i\ne j\).  Lemma~\ref{lem:pivot} gives elementary products
\(E,F\), represented by lists \(\ell,r\), such that
\[
  B=EAF,\qquad
  B_{jj}=1,\qquad B_{ij}=B_{ji}=0,
\]
with
\[
  \abs\ell\le2,\qquad \abs r\le3.
\]

For every \(k\ne j\), right multiply by
\[
  e_{jk}(-B_{jk}).
\]
These transvections have a common source column and distinct target
columns, so they do not interfere with one another.  If \(T\) is
their product and \(C=BT\), then
\[
  C_{jc}=\begin{cases}1,&c=j,\\0,&c\ne j.\end{cases}
\]
The list \(t\) of right-clearing factors has length at most \(n\).

For every \(k\ne j\), left multiply by
\[
  e_{kj}(-C_{kj}).
\]
Again the operations have distinct target rows and do not interfere.
If \(S\) is their product and \(D=SC\), then
\[
  D_{rc}=
  \begin{cases}
    1,&r=c=j,\\
    0,&r=j,\ c\ne j,\\
    0,&r\ne j,\ c=j,\\
    C_{rc},&r,c\ne j.
  \end{cases}
\]
Hence \(D\in\GL_{n-1}^{(j)}(R)\).  The list \(s\) of
left-clearing factors also has length at most \(n\).

Solving \(D=SEAF T\) for \(A\) gives
\[
  A=E^{-1}S^{-1}DT^{-1}F^{-1}.
\]
At the word level, inverse products are represented by reversing
the factor lists and inverting every transvection.  The inverse of
\(e_{ab}(z)\) is \(e_{ab}(-z)\), so all these factors remain
elementary.  The resulting word has the exact form
\[
  \ell^{-1}_{\mathrm{rev}}\,
  s^{-1}_{\mathrm{rev}}\,
  [D]\,
  t^{-1}_{\mathrm{rev}}\,
  r^{-1}_{\mathrm{rev}},
\]
and therefore length at most
\[
  2+n+1+n+3=2n+6.
\]
\end{proof}

\begin{remark}
The fixed coordinate \(j\) is chosen before \(A\).  Allowing the
coordinate block to depend on \(A\) gives a different, weaker width
statement; the fixed-coordinate form is what feeds the uniform
quasi-cocycle globalization below.
\end{remark}

\section{From relative control to global rigidity}\label{sec:globalization}

Let \(R\) be a finite-type \(\F_2\)-algebra.  Choose a finite set
\(X\) and a surjection
\[
  \F_2\langle X\rangle\twoheadrightarrow R.
\]
Elementary matrices are functorial in the coefficient ring, so the
root estimate of Theorem~\ref{thm:root} descends to \(R\).  Coordinate
permutations then give the same estimate for every root subgroup in
rank four.

Two larger pieces are used in the assembly: a last-column subgroup
and a last-row subgroup.  The quasi-cocycle is bounded on both.
The stabilized rank-three coordinate group normalizes these pieces.
Kazhdan control for the rank-three group converts boundedness on the
normalized root system into boundedness on the coordinate group,
provided the ambient representation has no invariant vectors.  This
is the only place where the quotient by the trivial representation
enters.

We record the relative-to-global step explicitly.  For a subset
\(S\subseteq G\), say that quasi-cocycles are \emph{bounded on
\(S\)} if every quasi-cocycle into a unitary representation without
invariant vectors has bounded restriction to \(S\).

\begin{lemma}[Kazhdan globalization on a normalized set]
\label{lem:kazhdan-globalization}
Let \(G\) have property~\((T)\), let \(S\subseteq G\) generate
\(G\), and let \(H\le G\).  Suppose
\[
  h^{-1}Sh\subseteq S\qquad(h\in H)
\]
and quasi-cocycles are bounded on \(S\).  Then \((G,H)\) has
relative property \((TT)/T\).
\leanverified{PropertyTT/PaperStatements}{PropertyTTPaper.kazhdan_normalizedSet_globalization}
\end{lemma}

\begin{proof}
Fix a unitary representation \(\pi\) of \(G\) without nonzero
invariant vectors and a quasi-cocycle \(b\) of defect at most \(D\).
Choose a Kazhdan pair \((K,\varepsilon)\), and choose once and for
all an \(S\cup S^{-1}\)-word for each \(k\in K\).  Let \(L\) be
the sum of their lengths.  If \(\norm{b(s)}\le C\) for \(s\in S\),
comparison of the quasi-cocycle identities for \(sh\) and
\(h(h^{-1}sh)\) gives
\[
  \norm{\pi(s)b(h)-b(h)}\le 2C+2D
  \qquad(s\in S\cup S^{-1},\ h\in H).
\]
The displacement estimate for a product of unitaries therefore gives
\[
  \norm{\pi(k)b(h)-b(h)}
  \le L(2C+2D)\qquad(k\in K).
\]
Since \(\pi\) has no invariant vectors, the Kazhdan inequality
supplies some \(k\in K\) with
\(
  \varepsilon\norm{b(h)}
  \le\norm{\pi(k)b(h)-b(h)}.
\)
Thus \(\norm{b(h)}\le L(2C+2D)/\varepsilon\), uniformly in
\(h\in H\).
\end{proof}

\begin{lemma}[Leavitt rank equivalence]
\label{lem:rank-equivalence}
If \(R\) contains a binary Leavitt family, then for all
\(m,n\ge2\) there is a group isomorphism
\[
  \EL_m(R)\cong\EL_n(R).
\]
\leanverified{PropertyTT/PaperStatements}{PropertyTTPaper.leavitt_elementaryRankEquivalence}
\end{lemma}

\begin{proof}
Iterating the binary prefix code gives a ring isomorphism
\(M_q(R)\cong R\) for every \(q\ge2\).  Functoriality in the
coefficient ring and block flattening give
\[
 \EL_m(R)
 \cong \EL_m(M_n(R))
 \cong \EL_{mn}(R)
 \cong \EL_n(M_m(R))
 \cong \EL_n(R).
\]
All maps in this chain preserve elementary root generators, so the
display restricts from general linear groups to the stated
elementary groups.
\end{proof}

\begin{lemma}[Elementary equals general linear]\label{lem:gl-el}
Assume that \(R\) has a binary Leavitt family, satisfies~\eqref{eq:sandwich},
and satisfies
\(\diag(u,1)\in\EL_2(R)\) for every \(u\in R^\times\).  Then
\[
  \EL_n(R)=\GL_n(R)\qquad(n\ge2).
\]
\leanverified{PropertyTT/PaperStatements}{PropertyTTPaper.elementaryGroup_eq_generalLinear}
\end{lemma}

\begin{proof}
The single-sandwich elimination reduces an arbitrary element of
\(\GL_2(R)\), by elementary row and column operations, to
\(\diag(u,1)\) for some \(u\in R^\times\).  The diagonal hypothesis
therefore gives \(\GL_2(R)=\EL_2(R)\).

Apply this rank-two equality over the coefficient ring \(M_n(R)\).
Under block flattening, an elementary transvection over \(M_n(R)\)
splits into scalar transvections, and hence
\[
  \EL_2(M_n(R))=\GL_2(M_n(R))
  \quad\Longrightarrow\quad
  \EL_{2n}(R)=\GL_{2n}(R).
\]
Equivalently, one may write the same flattening after interchanging
the two finite index factors as
\[
  \EL_n(M_2(R))=\GL_n(M_2(R)).
\]
The binary Leavitt family supplies a coefficient isomorphism
\(M_2(R)\cong R\), and transport along it gives
\(\EL_n(R)=\GL_n(R)\).  This argument uses block flattening only in
the form in which elementary block transvections decompose into
ordinary elementary transvections.
\end{proof}

\begin{proposition}[Rank-four assembly]\label{prop:rank-four}
Under the hypotheses of Theorem~\ref{thm:B}, \(\EL_4(R)\) has
property \((TT)/T\).
\leanverified{PropertyTT/PaperStatements}{PropertyTTPaper.finitePresentation_rankFour_hasTTmodT}
\end{proposition}

\begin{proof}
Choose a finite free-algebra presentation of \(R\).  Theorem~\ref{thm:root}
descends through this quotient and gives relative
\((TT)/T\) for \(X_{03}(R)\).  Coordinate permutations give it for
the six roots
\[
  X_{03},X_{13},X_{23},X_{30},X_{31},X_{32}.
\]

The assembly has the following dependency chain:
\[
 \boxed{
 \begin{gathered}
   \text{six root bounds}
   \Longrightarrow \text{last row and last column are bounded}\\
   \Longrightarrow \text{the rank-three core is bounded}
   \Longrightarrow \text{coordinate blocks are bounded}\\
   \xLongrightarrow{\text{Theorem A}}
   \text{the whole rank-four group is bounded}.
 \end{gathered}}
\]
Here are the details.  Put
\[
  C=X_{03}X_{13}X_{23},
  \qquad C^-=X_{30}X_{31}X_{32}.
\]
The three roots in each product commute, so the quasi-cocycle
identity makes \(b\) bounded on \(C\) and on \(C^-\).  The set
\(S=C\cup C^-\) generates \(\EL_4(R)\): roots meeting the last
coordinate lie in \(S\), and every remaining root is a Steinberg
commutator of one last-column and one last-row root.

Property~\((T)\) for \(\EL_3(R)\) follows from
Ershov--Jaikin-Zapirain; equivalently, in the formal proof it follows
by descending the explicit finite Kazhdan system for the free
algebra.  Lemma~\ref{lem:rank-equivalence} transports property
\((T)\) to \(\EL_4(R)\).  The stabilized rank-three subgroup
\(H=\EL_3(R)\) normalizes both \(C\) and \(C^-\), hence normalizes
\(S\).  Lemma~\ref{lem:kazhdan-globalization} now proves that
\((\EL_4(R),H)\) has relative property \((TT)/T\).  Notice that the
no-invariant-vectors hypothesis is used precisely in the Kazhdan
inequality inside that lemma.

By Lemma~\ref{lem:gl-el}, the elementary groups in ranks three and
four are the full general linear groups.  Under these
identifications the stabilized core is exactly the coordinate block
\(\GL_3^{(3)}(R)\).  Finally, coordinate permutations carry this
relative statement to \(\GL_3^{(j)}(R)\) for every
\(j\in\{0,1,2,3\}\).

Each elementary transvection lies in a coordinate block avoiding its
two indices.  Hence the union of one fixed coordinate block and all
root subgroups is bounded for the quasi-cocycle.  Theorem~\ref{thm:A}
writes every element of \(\GL_4(R)\) as a product of at most \(14\)
elements of that union.  Lemma~\ref{lem:products} below, or directly
the corresponding bounded-product lemma, makes the quasi-cocycle
bounded on \(\GL_4(R)\).  Transporting through
\(\EL_4(R)=\GL_4(R)\) proves the proposition.
\end{proof}

The passage from a bounded generating subset to the whole group
uses the following elementary estimate.

\begin{lemma}[Products of controlled pieces]\label{lem:products}
Let \(b\) be a quasi-cocycle of defect at most \(D\), and suppose
\(\norm{b(g_i)}\le C\) for a list \((g_1,\ldots,g_m)\).  Then
\[
  \norm{b(g_1\cdots g_m)}
  \le mC+(m+1)D.
\]
\leanverified{PropertyTT/PaperStatements}{PropertyTTPaper.quasiCocycle_list_product_bound}
\end{lemma}

\begin{proof}
Induct on the list.  The empty-list case is
\(\norm{b(1)}\le D\).  For a nonempty list, apply~\eqref{eq:quasicocycle}
to the head and the product of the tail, use
that the representation is isometric, and invoke the inductive
bound.  The displayed estimate is a slightly loose uniform form that
also covers the empty list.
\end{proof}

\begin{proof}[Proof of Theorem~\ref{thm:B}]
Proposition~\ref{prop:rank-four} gives the result in rank four.  For
arbitrary \(n\ge3\), Lemma~\ref{lem:rank-equivalence} supplies an
explicit group equivalence \(\EL_n(R)\cong\EL_4(R)\).
Property \((TT)/T\) is invariant under group equivalence, so the
rank-four conclusion transports to rank \(n\).
\end{proof}

\section{The binary Leavitt specialization}\label{sec:leavitt}

Let \(L\) be the binary Leavitt algebra~\cite{Leavitt},
\[
  L=L_{\F_2}(1,2)
  =\F_2\langle x_0,x_1,y_0,y_1\rangle/
  (y_i x_j-\delta_{ij},\,x_0y_0+x_1y_1-1).
\]
The defining generators themselves form a binary Leavitt family,
and the presentation makes \(L\) a finite-type \(\F_2\)-algebra.
Ara--Goodearl--Pardo prove that \(L\) is purely infinite simple
\cite[Theorem~4.2]{AraGoodearlPardo}; their
Theorem~1.6 then gives the single-sandwich property.  It remains to
discuss the unstable diagonal condition.

\begin{proposition}[Unstable diagonal reduction]\label{prop:diag}
For every \(u\in L^\times\),
\[
  \diag(u,1)\in\EL_2(L).
\]
Consequently \(\GL_2(L)=\EL_2(L)\).
\leanverified{KOne/PaperStatements}{KOnePaper.diagUnit_mem_elementary}
\leanverified{KOne/PaperStatements}{KOnePaper.elementaryGroup_two_eq_top}
\end{proposition}

\begin{proof}
Ara--Brustenga--Corti\~nas give, for a regular coefficient ring, the
exact sequence determined by \(1-N_E^t\)~\cite{AraBrustengaCortinas}.
For the graph with one vertex and two loops this map is multiplication
by \(-1\), hence
\[
  K_i(L)=0
\]
for every \(i\in\mathbb Z\), in particular in degree one.

Ara--Goodearl--Pardo prove that a unital purely infinite
simple ring is a GE-ring and that
\[
  K_1(R)\cong R^\times_{\mathrm{ab}}
\]
\cite[Theorem~2.4]{AraGoodearlPardo}.  Thus \(K_1(L)=0\) makes each
unit a product of commutators.  The standard \(2\times2\) Whitehead
identity places \(\diag([a,b],1)\) in \(\EL_2(L)\); the GE property
then gives \(\GL_2(L)=\EL_2(L)\), proving
Proposition~\ref{prop:diag}.
The endpoint is therefore established by the cited structural
theorems, which prove Proposition~\ref{prop:diag}.
\end{proof}

\begin{proof}[Proof of Theorem~\ref{thm:C}]
The binary Leavitt algebra is a nontrivial finite-type
\(\F_2\)-algebra, its canonical generators form a binary Leavitt
family, pure infiniteness supplies~\eqref{eq:sandwich}, and
Proposition~\ref{prop:diag} supplies the elementary diagonal
condition.  Theorem~\ref{thm:B} applies.
\end{proof}

\section{Further directions}\label{sec:directions}

The characteristic-two hypothesis enters most visibly
in~\eqref{eq:fixed}, where every root element has order two.  Over a
field of characteristic \(p\), the corresponding root has order
\(p\), and iterating the quasi-cocycle identity gives
\[
  b(1)\approx
  b(r)+\pi(r)b(r)+\cdots+\pi(r)^{p-1}b(r).
\]
On the root-fixed space this is \(p\) times the fixed projection of
\(b(r)\), with error at most \((p-1)D\).  Thus the fixed-space
argument has a plausible finite-characteristic analogue.  Extending
the theorem requires uniform replacements for the
characteristic-two rank-three control and the sign-free coordinate
symmetries.

The algebraic factorization theorem suggests a separate question.
For which natural classes of simple or purely infinite rings does
the single-sandwich condition hold, and when can the coordinate
block in Theorem~\ref{thm:A} itself be controlled by elementary
subgroups?  A satisfactory answer would extend the present
synthesis from the Leavitt setting to a broader rigidity criterion.

\subsection*{Formal companion}

\ifanonymous
An accompanying Lean formalization is available in the supplementary
material.
\else
An accompanying Lean formalization is available
\href{\leanrepo}{online}.
\fi

\begin{thebibliography}{99}

\bibitem{AraBrustengaCortinas}
P.~Ara, M.~Brustenga, and G.~Corti\~nas,
\emph{\(K\)-theory of Leavitt path algebras},
M\"unster~J.\ Math. \textbf{2} (2009), 5--34.

\bibitem{AraGoodearlPardo}
P.~Ara, K.~R. Goodearl, and E.~Pardo,
\emph{\(K_0\) of purely infinite simple regular rings},
K-Theory \textbf{26} (2002), no.~1, 69--100.

\bibitem{ErshovJaikin}
M.~Ershov and A.~Jaikin-Zapirain,
\emph{Property \((T)\) for noncommutative universal lattices},
Invent.\ Math. \textbf{179} (2010), no.~2, 303--347.

\bibitem{Leavitt}
W.~G. Leavitt,
\emph{The module type of a ring},
Trans.\ Amer.\ Math.\ Soc. \textbf{103} (1962), 113--130.

\bibitem{MimuraTT}
M.~Mimura,
\emph{Property \((TT)\) modulo \(T\) and homomorphism
superrigidity into mapping class groups},
preprint, 2011,
\href{https://arxiv.org/abs/1106.3769}{arXiv:1106.3769}.

\bibitem{MimuraThesis}
M.~Mimura,
\emph{Rigidity theorems for universal and symplectic universal lattices},
doctoral thesis, University of Tokyo, 2011,
\href{https://doi.org/10.15083/00004046}{doi:10.15083/00004046}.

\bibitem{OpenAITenProofs}
OpenAI,
\emph{Ten advances in mathematics and theoretical computer science},
2026, Chapter~3: \emph{Nonsofic groups exist}, pp.~78--95,
\href{https://cdn.openai.com/pdf/ten-proofs-oai.pdf}{proof volume};
see also the
\href{https://openai.com/index/ten-advances-in-mathematics/}{project announcement}.

\end{thebibliography}

\end{document}
